% last updated in April 2002 by Antje Endemann
% Based on CVPR 07 and LNCS, with modifications by DAF, AZ and elle, 2008 and AA, 2010, and CC, 2011; TT, 2014; AAS, 2016

\documentclass[runningheads]{llncs}
\usepackage{graphicx}
\usepackage{amsmath,amssymb} % define this before the line numbering.
\usepackage{ruler}
\usepackage{multirow}
\usepackage{booktabs}
\usepackage{threeparttable}

\usepackage{pifont}
\usepackage{color}
\usepackage[width=122mm,left=12mm,paperwidth=146mm,height=193mm,top=12mm,paperheight=217mm]{geometry}
\usepackage[pagebackref,breaklinks,colorlinks,citecolor=blue]{hyperref}

\usepackage{longtable}
\usepackage{booktabs}
\usepackage{multirow}
\usepackage{xltabular}
\usepackage{subcaption} % 用于并排显示子图
\usepackage{makecell}

\begin{document}
% \renewcommand\thelinenumber{\color[rgb]{0.2,0.5,0.8}\normalfont\sffamily\scriptsize\arabic{linenumber}\color[rgb]{0,0,0}}
% \renewcommand\makeLineNumber {\hss\thelinenumber\ \hspace{6mm} \rlap{\hskip\textwidth\ \hspace{6.5mm}\thelinenumber}}
% \linenumbers
\pagestyle{headings}
\mainmatter
\def\ECCV16SubNumber{13487}  % Insert your submission number here
\newcommand{\methodname}{TextComp} 
\title{Appendix} % Replace with your title

\titlerunning{ECCV-26 submission ID \ECCV16SubNumber}

\authorrunning{ECCV-26 submission ID \ECCV16SubNumber}

\author{Anonymous ECCV submission}
\institute{Paper ID \ECCV16SubNumber}


\maketitle

\section{A. Attribute Word Generation Details}
\label{sec:appendix_attribute}

We leverage several foundation language models, including Qwen and Llama, to generate attribute words for action categories. Due to the inherent randomness in the generation process and varying model architectures, the produced attribute words can differ across models. To ensure consistency and quality of the generated attributes, we design a unified template instruction and apply it uniformly to all foundation models. The general instruction prompt is formulated as follows:

\begin{quote}
\textit{Describe the visual characteristics of the action ``\{class\_name\}'' from the \{dataset\_name\} video action recognition dataset. Provide \{n\_words\} single descriptive words (NOT phrases) about the motion, posture, and movement patterns. Format your response as follows: word word word...}
\end{quote}

We incorporate several explicit constraints in the prompt design to ensure output quality. First, we emphasize generating single words rather than phrases to maintain semantic granularity and facilitate subsequent processing. Second, we explicitly instruct the models to generate English words only, as multilingual models may occasionally produce outputs in other languages or even non-verbal symbols such as emojis. Third, we specify a structured output format (space-separated words) to enable seamless data loading in downstream processes without requiring complex parsing.

After generation, we apply a post-processing pipeline to filter the raw outputs: (1) removing non-alphabetic characters and symbols, (2) eliminating duplicate words within each category, and (3) discarding overly generic terms that appear across most action classes. This cleaning procedure ensures that the final attribute set contains discriminative and action-specific descriptors.

Table~\ref{tab:attribute_samples} presents sample attribute words for representative action categories from UCF101, HMDB51, Kinetics-400, and Something-Something v2 (SSv2) generated by different foundation models. We observe that while different models produce varied attributes, they consistently capture key semantic aspects of each action. The complete attribute word sets for all datasets and the generation code will be made publicly available upon publication.

% 注意：开头改成了 xltabular，并增加了 {\textwidth} 参数，第三列改成了 X
\begin{xltabular}{\textwidth}{l|l|X}
\caption{Sample attribute words generated by different foundation models for representative action categories across datasets.}
\label{tab:attribute_samples}\\
\toprule
\textbf{Dataset} & \textbf{Action} & \textbf{Sample Attributes} \\
\midrule
\endfirsthead

\multicolumn{3}{c}{\tablename\ \thetable\ -- \textit{Continued from previous page}} \\
\toprule
\textbf{Dataset} & \textbf{Action} & \textbf{Sample Attributes} \\
\midrule
\endhead

\midrule
\multicolumn{3}{r}{\textit{Continued on next page}} \\
\endfoot

\bottomrule
\endlastfoot

\multirow{4}{*}{UCF101} 
& Archery & \textbf{Qwen3:8b:} word shooting target arrow bow release aiming draw string nock fletching positioning stance balance focus concentration precision accuracy steady controlled relaxed tense trajectory flight path arc curved straight launch propulsion hit miss \\
& & \textbf{llama3.1:8b:} Standing Holding Steady Straight Forward Arm Backward Focused Calm Stable Still Aiming Tense Relaxed Vertical Horizontal Level Balanced Firm Grip Strong Controlled Smooth Deliberate Precision Concentrated Calibrated Stance Firmly Secure Alert \\
\cmidrule{2-3}
& Juggling Balls & \textbf{Qwen3:8b:} ball juggling juggler dynamic rhythmic fluid fast quick controlled precision continuous bouncing balancing dexterous coordinated repetitive circular spinning rotating rapid swift trick steady precise skillful motion action performance movement timing challenge \\
& & \textbf{llama3.1:8b:} Multiple Moving Fast Upward Forward Rotating Circular Smooth Continuous Alternating Quick Frequent Hand Arm Ball Catching Dropping Repeating Synchronized Entertaining Energetic Active Dynamic Flicking Throwing Intercepting Swinging Bouncing Rising Falling Oscillating Agile \\
\midrule

\multirow{4}{*}{HMDB51}
& brush hair & \textbf{Qwen3:8b:} brush comb style arrange neat tidy smooth glossy tangle untangle curl straighten fix secure manipulate direct position align adjust prepare maintain clean detangle brushing combing hair pin sleek polished groom manage \\
& & \textbf{llama3.1:8b:} Moving Slow Gentle Small Upward Forward Smooth Soft Side Curved Twisting Waving Horizontal Vertical Fast Quick Light Flicking Backward Circular Whirling Rotating Turning Swaying Swinging Pushing Pressing Steady Controlled Precise \\
\cmidrule{2-3}
& sword exercise & \textbf{Qwen3:8b:} blade slash parry lunge feint dodge counter crouch attack block spin twirl thrust recover balance stance stab deflect hook guard camber recoil drive grip twist draw sheath vault pommel edge hilt motion \\
& & \textbf{llama3.1:8b:} Twisting Raising Lowering Arched Forward Bending Turning Swaying Rotating Lifting Straight Holding Moving Linear Steady Vertical Horizontal Sharp Quick Controlled Flexible Fluid Dynamic Strong Stable Level Upward Downward Swinging Angular Balanced Smooth \\
\midrule

\multirow{4}{*}{Kinetics-400}
& air drumming & \textbf{Qwen3:8b:} rhythmic energetic dynamic fluid expressive creative imaginative artistic vibrant powerful dramatic graceful agile confident commanding passionate spontaneous cathartic synchronized innovative dexterous versatile playful frenetic epic captivating stylized syncopated effortless controlled engaging performative \\
& & \textbf{llama3.1:8b:} Moving Rhythmic Flicking Quick Light Agile Twirling Bent Relaxed Upward Lateral Alternating Smooth Fast Precise Arm Waist Rotating Hip Swiveling Calm Expressive Dynamic Playful Bouncy Energetic Free-flowing Unpredictable Whirling Loose \\
\cmidrule{2-3}
& balloon blowing & \textbf{Qwen3:8b:} inflating blowing puffing expanding swelling stretching ballooning airflow pressure tension elasticity movement action motion dynamic fluid velocity force resistance control coordination hand mouth breathing exhalation intake release balloon air \\
& & \textbf{llama3.1:8b:} Inhaling Moving Upward Forward Bending Twisting Rotating Smooth Gentle Slow Steady Focused Concentrated Inward Curving Circular Rhythmic Repetitive Consistent Precise Controlled Intense Concentric Spiral Ascending Gradual Linear Horizontal Slightly Unstable Tensing \\
\midrule

\multirow{4}{*}{SSv2}
& Closing something & \textbf{Qwen3:8b:} close shut fasten seal lock tighten secure bolt clamp plug cover wrap bind tie zip button latch click slide turn rotate twist screw unscrew press push \\
& & \textbf{llama3.1:8b:} Focusing Concentrating Grasping Tightening Pressing Moving Approaching Reaching Touching Compressing Squeezing Applying Slowing Stopping Folding Tucking Securing Holding Resting Laying Placing Inserting Aligning Positioning Adjusting Snapping Clicking Locking Fastening Secure \\
\cmidrule{2-3}
& Picking something up & \textbf{Qwen3:8b:} grasping lifting retrieving uplifting handling manipulating acquiring obtaining securing gathering collecting accumulating amassing fetching reaching snatching grabbing plucking picking selecting choosing taking seizing capturing extracting pulling hoisting uprooting transferring transporting relocating \\
& & \textbf{llama3.1:8b:} Reaching Bending Forward Upward Smooth Slow Deliberate Controlled Precise Arm Hand Open Fingers Spread Grasping Pulling Outward Motion Linear Horizontal Tilting Forwardly Quick Rigid Stable Balanced Steady Level Lifting Rising Gradual Purposeful \\

% 结尾也改成了 xltabular
\end{xltabular}


\section{B. Semantic Discriminability and Prompt Robustness Analysis}
\label{sec:semantic_discriminability}

To further investigate the impact of the \texttt{\{dataset\_name\}} prompt variable and demonstrate our method's in-the-wild generalization capabilities, we conduct a quantitative semantic separability analysis in the feature space. The core objective is to show that our proposed Attribute Enhancement mechanism effectively increases the semantic distance (i.e., discriminability) between different action categories, even when applied to unknown, open-world actions without dataset-specific guidance.



\begin{table}[t]
\centering
\caption{Quantitative comparison of inter-class semantic similarity. Lower off-diagonal mean and top-1 confusion values indicate higher discriminative power in the feature space}
\label{tab:semantic_sim}
\begin{tabular}{lcc}
\hline
Strategy & Off-diagonal Mean Sim. $\downarrow$ & Top-1 Confusion Mean $\downarrow$ \\
\hline
        Base & 0.8418 & 0.9087 \\
        w/o Dataset Name & 0.8018 & 0.8926 \\
        w/ Dataset Name & \textbf{0.7593} & \textbf{0.8604} \\
\hline
\end{tabular}
\end{table}


\begin{figure*}[t]
    \centering
    % 请确保图片路径与你本地保存的路径一致
    \includegraphics[width=0.32\linewidth]{figure/heatmap/heatmap_base.png}
    \hfill
    \includegraphics[width=0.32\linewidth]{figure/heatmap/heatmap_ours2.png}
    \hfill
    \includegraphics[width=0.32\linewidth]{figure/heatmap/heatmap_ours1.png}
    \caption{Inter-class cosine similarity heatmaps of CLIP text features. \textbf{Left:} Base (Pure Class Name). \textbf{Middle:} Ours (w/o Dataset Name). \textbf{Right:} Ours (w/ Dataset Name). Lighter/cooler colors in the off-diagonal areas represent lower similarity and thus better inter-class separability.}
    \label{fig:heatmaps}
\end{figure*}


While foundation models exhibit inherent stochasticity in text generation, the underlying semantic distribution of the generated attributes should remain robust. To verify this, we extract text features using the pre-trained CLIP text encoder for three distinct settings on the HMDB-51 action vocabulary:
\begin{enumerate}
    \item \textbf{Base:} Using only the raw action class names without any attribute enhancement.
    \item \textbf{w/o Dataset Name:} Using attributes generated completely without the dataset context, simulating an in-the-wild scenario.
    \item \textbf{w/ Dataset Name:} Using our full prompt including the dataset name, identical to the main paper's setting.
\end{enumerate}

We compute the inter-class cosine similarity matrix for all categories. A lower off-diagonal mean similarity indicates a greater distance between different categories in the feature space, reflecting higher semantic discriminative power. The quantitative results and corresponding heatmaps are presented in Table~\ref{tab:semantic_sim} and Figure~\ref{fig:heatmaps}.


\textbf{Analysis:} As shown in Table~\ref{tab:semantic_sim} and Figure~\ref{fig:heatmaps}, relying solely on raw class names (Base) results in a high inter-class similarity ($0.8418$). Applying our Attribute Enhancement \textit{without} the dataset name reduces this similarity to $0.8018$. This confirms that our generated attributes successfully extract discriminative semantic features and improve inter-class separability, demonstrating robust in-the-wild generalization even without dataset-specific priors. Finally, including the dataset name yields the highest discriminability (the lowest similarity at $0.7593$). This validates our prompt design for academic benchmarks, as dataset-specific context further aids foundation models in domain-level semantic disambiguation. This feature-space behavior perfectly aligns with the prompt ablation study in our main text.


\section{C. Broader Applicability: Zero-shot Video-Text Retrieval}
\label{sec:retrieval_applicability}

To demonstrate the generalizability of our method, we evaluate its representation power on a broader vision-language task: \textbf{Zero-shot Video-Text Retrieval}.

Unlike zero-shot classification, retrieval searches through a large gallery of unconstrained videos given a textual query (Text-to-Video), or vice versa (Video-to-Text). In practical scenarios, we formulate our generated motion and posture attributes as a zero-shot \textit{Query Expansion} module. Our objective is to determine whether augmenting a basic category query with our generated fine-grained attributes produces a stronger, more robust text representation capable of aligning with visual details more accurately.

\textbf{Experimental Protocol:}
We utilize the HMDB-51 test set for this evaluation.
The text queries are formulated in two ways:
\begin{enumerate}
    \item \textbf{Base:} \textit{``a video of \{class\_name\}''}
    \item \textbf{Ours:} \textit{``a video of \{class\_name\}. attributes: \{generated\_attributes\}''}
\end{enumerate}
We extract features using the pre-trained CLIP model and compute the cosine similarity
between the text queries and all video features.
Performance is measured using standard retrieval metrics:
Recall@K (R@1, R@5, R@10) and Median Rank (MdR).




\begin{table}[h]
\centering
\caption{Zero-shot Video-Text Retrieval on HMDB-51.}
\label{tab:retrieval_results}
\begin{tabular}{lcccccccc}
\hline
 & \multicolumn{4}{c}{Text-to-Video} & \multicolumn{4}{c}{Video-to-Text} \\
Method & R@1 & R@5 & R@10 & MdR & R@1 & R@5 & R@10 & MdR \\
\hline
Base  & 49.02 & 66.67 & 78.43 & 2.0 & 23.33 & 50.78 & 63.01 & 5.0 \\
Ours  & \textbf{50.98} & \textbf{72.55} & \textbf{82.35} & \textbf{1.0} & \textbf{23.73} & \textbf{50.85} & \textbf{63.27} & \textbf{5.0} \\
\hline
\end{tabular}
\end{table}


\textbf{Analysis:} As reported in Table~\ref{tab:retrieval_results}, leveraging our Attribute Enhancement for query expansion consistently improves both Text-to-Video and Video-to-Text retrieval performance. Notably, in Text-to-Video retrieval, our method significantly boosts R@5 from 66.67\% to 72.55\% (+5.88\%) and improves the Median Rank from 2.0 to 1.0. This compellingly demonstrates that our method is not merely a task-specific classification heuristic. By decomposing macro-actions into dense, motion-informed descriptors, it builds a richer semantic intermediary space.
\clearpage


\end{document}
